\section{Dimer models and perfect matching polygons}\label{sec_pre_dimer}

\subsection{What is a dimer model?}
We first introduce dimer models. Some ideas and concepts contained in this subsection are originally derived from theoretical physics (e.g., \cite{FHK,HV}).

A \emph{dimer model} (or \emph{brane tiling}) $\Gamma$ is a finite bipartite graph on the real two-torus~$\TT\coloneqq\RR^2/\ZZ^2$; that is, the set~$\Gamma_0$ of nodes is divided into two parts~$\Gamma_0^+$,~$\Gamma_0^-$, and the set~$\Gamma_1$ of edges consists of the edges connecting nodes in~$\Gamma_0^+$ with those in~$\Gamma_0^-$.
In order to make the situation clear, we color the nodes in~$\Gamma_0^+$ white, and those in~$\Gamma_0^-$ black.
A connected component of $\TT{\setminus}\Gamma_1$ is called a \emph{face} of $\Gamma$, and we denote by~$\Gamma_2$ the set of faces.
We also obtain the bipartite graph $\widetilde{\Gamma}$ on $\RR^2$ induced via the universal cover $\RR^2\rightarrow\TT$.
We call $\widetilde{\Gamma}$ the universal cover of the dimer model~$\Gamma$.
For example, the bipartite graph shown on the left side of Figure~\ref{ex_dimer_4b} is a dimer model
where the outer frame is the fundamental domain of~$\TT$.

As the dual of a dimer model $\Gamma$, we define the quiver $Q_\Gamma$ associated with $\Gamma$.
Namely, we assign a vertex dual to each face in $\Gamma_2$, and an arrow dual to each edge in $\Gamma_1$.
The orientation of arrows is determined so that the white node is on the right of the arrow.
For example, the right side of Figure~\ref{ex_dimer_4b} is the quiver associated with the dimer model on the left.
Sometimes we simply denote the quiver $Q_\Gamma$ by $Q$.

%%%%%%%%%%%%%%%%%%%%%
%%% type 4b %%%%%%%%%%%%%
%%%%%%%%%%%%%%%%%%%%%
\newcommand{\basicdimerB}{
%coordinate setting
\coordinate (W1) at (5,1); \coordinate (W2) at (1,3); \coordinate (W3) at (3,5);
\coordinate (B1) at (3,1); \coordinate (B2) at (1,5); \coordinate (B3) at (5,3);
\draw[line width=\edgewidth] (0,0) rectangle (6,6);
%edge
\draw[line width=\edgewidth] (W2)--(B1)--(W1)--(B3)--(W3)--(B2)--(W2); \draw[line width=\edgewidth] (B1)--(W3);
\draw[line width=\edgewidth] (0,3)--(W2); \draw[line width=\edgewidth] (3,0)--(B1); \draw[line width=\edgewidth] (6,0)--(W1);
\draw[line width=\edgewidth] (6,3)--(B3); \draw[line width=\edgewidth] (3,6)--(W3); \draw[line width=\edgewidth] (0,6)--(B2);
%black
\draw [line width=\nodewidth, fill=black] (B1) circle [radius=\noderad] ;
\draw [line width=\nodewidth, fill=black] (B2) circle [radius=\noderad] ;
\draw [line width=\nodewidth, fill=black] (B3) circle [radius=\noderad] ;
%white
\draw [line width=\nodewidth, fill=white] (W1) circle [radius=\noderad] ;
\draw [line width=\nodewidth, fill=white] (W2) circle [radius=\noderad] ;
\draw [line width=\nodewidth, fill=white] (W3) circle [radius=\noderad] ;
}

\begin{figure}[h!]\centering
\newcommand{\edgewidth}{0.07cm}
\newcommand{\nodewidth}{0.07cm}
\newcommand{\noderad}{0.24} %radius_of_node
\newcommand{\arrowwidth}{0.08cm}
\begin{tikzpicture}
\node (DM) at (0,0)
{\scalebox{0.45}{
\begin{tikzpicture}
\basicdimerB
\end{tikzpicture}
} };

\node (QV) at (5,0)
{\scalebox{0.45}{
\begin{tikzpicture}[sarrow/.style={black, -latex, very thick}]
%vertex
\node (Q1) at (2,3.5){{\LARGE$1$}}; \node (Q0) at (4,2.5){{\LARGE$0$}};
\node (Q3) at (5,5){{\LARGE$3$}}; \node (Q2) at (1,1){{\LARGE$2$}};
\draw[line width=\edgewidth] (0,0) rectangle (6,6);

%edge
\draw[lightgray, line width=\edgewidth] (W2)--(B1)--(W1)--(B3)--(W3)--(B2)--(W2); \draw[lightgray, line width=\edgewidth] (B1)--(W3);
\draw[lightgray, line width=\edgewidth] (0,3)--(W2); \draw[lightgray, line width=\edgewidth] (3,0)--(B1);
\draw[lightgray, line width=\edgewidth] (6,0)--(W1); \draw[lightgray, line width=\edgewidth] (6,3)--(B3);
\draw[lightgray, line width=\edgewidth] (3,6)--(W3); \draw[lightgray, line width=\edgewidth] (0,6)--(B2);

%black
\draw [line width=\nodewidth, draw=lightgray, fill=lightgray] (B1) circle [radius=\noderad] ;
\draw [line width=\nodewidth, draw=lightgray, fill=lightgray] (B2) circle [radius=\noderad] ;
\draw [line width=\nodewidth, draw=lightgray, fill=lightgray] (B3) circle [radius=\noderad] ;
%white
\draw [line width=\nodewidth, draw=lightgray, fill=white] (W1) circle [radius=\noderad] ;
\draw [line width=\nodewidth, draw=lightgray, fill=white] (W2) circle [radius=\noderad] ;
\draw [line width=\nodewidth, draw=lightgray, fill=white] (W3) circle [radius=\noderad] ;

\draw[sarrow, line width=\arrowwidth] (Q0)--(Q1); \draw[sarrow, line width=\arrowwidth] (Q3)--(Q0); \draw[sarrow, line width=\arrowwidth] (Q1)--(Q2);
\draw[sarrow, line width=\arrowwidth] (Q2)--(3,0); \draw[sarrow, line width=\arrowwidth] (3,6)--(Q3); \draw[sarrow, line width=\arrowwidth] (Q3)--(4.7,6);
\draw[sarrow, line width=\arrowwidth] (4.7,0)--(Q0); \draw[sarrow, line width=\arrowwidth] (Q0)--(6,1.5); \draw[sarrow, line width=\arrowwidth] (0,1.5)--(Q2);
\draw[sarrow, line width=\arrowwidth] (Q2)--(0,3); \draw[sarrow, line width=\arrowwidth] (6,3)--(Q3); \draw[sarrow, line width=\arrowwidth] (Q3)--(6,4.5);
\draw[sarrow, line width=\arrowwidth] (0,4.5)--(Q1);
\draw[sarrow, line width=\arrowwidth] (1.3,0)--(Q2); \draw[sarrow, line width=\arrowwidth] (Q1)--(1.3,6); \draw[sarrow, line width=\arrowwidth] (Q2)--(0,0);
\draw[sarrow, line width=\arrowwidth] (6,6)--(Q3);
\end{tikzpicture}
} };

\end{tikzpicture}
\caption{A dimer model and the associated quiver.}\label{ex_dimer_4b}
\end{figure}

The \emph{valency} of a node is the number of edges incident to that node.
We say that a~node on a~dimer model is \emph{$n$-valent} if its valency is $n$.
We then define several operations on a dimer model.
The \emph{join move} is the operation removing a $2$-valent node and joining the two distinct nodes connected to it as shown in Figure~\ref{split_join}.
Thus, using join moves we obtain a dimer model having no $2$-valent nodes.
We say that a dimer model is \emph{reduced} if it has no $2$-valent nodes.
We can see that the quiver associated with a reduced dimer model contains no $2$-cycles.
On the other hand, there is the operation called the \emph{split move}, which inserts a $2$-valent node (see Figure~\ref{split_join}).

We say that two reduced dimer models $\Gamma$, $\Gamma^\prime$ are \emph{isomorphic}, which is denoted by $\Gamma\cong\Gamma^\prime$, if their underlying cell decompositions of $\TT$ are homotopy-equivalent.

\begin{figure}[h!]\centering
{\scalebox{1}{
\begin{tikzpicture}

\node at (3.8,0.5) {join move} ; \node at (3.8,-0.5) {split move} ;
\draw[->, line width=0.03cm] (2.6,0.15)--(5,0.15);
\draw[<-, line width=0.03cm] (2.6,-0.15)--(5,-0.15);

\node (Dimer_a) at (0,0)
{\scalebox{0.7}{
\begin{tikzpicture}
%vertex
\coordinate (b1) at (-1,0); \coordinate (b2) at (1,0);
\coordinate (w1) at (-2,1); \coordinate (w2) at (-2.5,0); \coordinate (w3) at (-2,-1);
\coordinate (w4) at (0,0); \coordinate (w5) at (2,1); \coordinate (w6) at (2,-1);

%edge
\draw[line width=0.05cm] (w1)--(b1); \draw[line width=0.05cm] (w2)--(b1);
\draw[line width=0.05cm] (w3)--(b1); \draw[line width=0.05cm] (b1)--(w4); \draw[line width=0.05cm] (w4)--(b2);
\draw[line width=0.05cm] (w5)--(b2); \draw[line width=0.05cm] (w6)--(b2);
%black
\filldraw [line width=0.05cm, fill=black] (b1) circle [radius=0.16] ; \filldraw [line width=0.05cm, fill=black] (b2) circle [radius=0.16] ;
%white
\draw [line width=0.05cm, fill=white] (w4) circle [radius=0.16] ;
\end{tikzpicture} }} ;


\node (Dimer_b) at (7,0)
{\scalebox{0.7}{
\begin{tikzpicture}
%vertex
\coordinate (b1) at (0,0);
\coordinate (w1) at (-1,1); \coordinate (w2) at (-1.5,0); \coordinate (w3) at (-1,-1);
\coordinate (w5) at (1,1); \coordinate (w6) at (1,-1);
%edge
\draw[line width=0.05cm] (w1)--(b1); \draw[line width=0.05cm] (w2)--(b1);
\draw[line width=0.05cm] (w3)--(b1); \draw[line width=0.05cm] (w5)--(b1); \draw[line width=0.05cm] (w6)--(b1);
%black
\filldraw [line width=0.05cm, fill=black] (b1) circle [radius=0.16] ;
%white
\end{tikzpicture} }} ;

\end{tikzpicture}
}}

\caption{An example of the join and split move.}\label{split_join}
\end{figure}

We sometimes focus on the following type of dimer models.

\begin{Definition}\label{def_hexagonal_square}
We say that a dimer model $\Gamma$ is \emph{hexagonal} (resp.\ \emph{rectangular})
if any face of~$\Gamma$ is hexagon (resp.\ rectangle) and any node of $\Gamma$ is $3$-valent (resp.\ 4-valent).
In particular, a~hexagonal (resp.\ rectangular) dimer model is homotopy-equivalent to a dimer model whose faces are all regular hexagons (resp.\ squares).
\end{Definition}

\subsection{Perfect matchings and the perfect matching polygon}\label{subsection_PM}

Next, we assign a lattice polygon to each dimer model.
For this purpose, we will introduce the notion of perfect matchings, and we construct a polygon called the \emph{perfect matching polygon}.

\begin{Definition}A \emph{perfect matching} (or \emph{dimer configuration}) on a dimer model $\Gamma$ is a subset $\sfP$ of $\Gamma_1$ such that each node is
the end point of precisely one edge in $\sfP$.
\end{Definition}

In general, not every dimer model necessarily has a perfect matching.
In this paper, we will mainly discuss consistent dimer models, and such dimer models have perfect matchings.
Moreover, we can extend a perfect matching $\sfP$ to one on $\widetilde{\Gamma}$ via the universal cover $\RR^2\rightarrow\TT$.
We call this a \emph{perfect matching} on $\widetilde{\Gamma}$, and use the same notation $\sfP$.
For example, some perfect matchings on the dimer model given in Figure~\ref{ex_dimer_4b} are shown in Figure~\ref{pm_4b}.
(This dimer model has eight perfect matchings in total.)

%%% type 4b perfect matching%%%
\begin{figure}[h!]\centering
\newcommand{\edgewidth}{0.07cm}
\newcommand{\nodewidth}{0.07cm}
\newcommand{\noderad}{0.24} %radius_of_node

\newcommand{\pmwidth}{0.5cm}
\newcommand{\pmcolor}{red}

\begin{tikzpicture}
\node at (0,-1.6) {$\sfP_0$};
\node at (3.2,-1.6) {$\sfP_1$}; \node at (6.4,-1.6) {$\sfP_2$}; \node at (9.6,-1.6) {$\sfP_3$}; \node at (12.8,-1.6) {$\sfP_4$};

\node (PM0) at (0,0)
{\scalebox{0.35}{
\begin{tikzpicture}
%perfect matching
\draw[line width=\pmwidth, color=\pmcolor] (W1)--(B3); \draw[line width=\pmwidth, color=\pmcolor] (W3)--(B1);
\draw[line width=\pmwidth, color=\pmcolor] (W2)--(B2);
\basicdimerB
\end{tikzpicture} }};

\node (PM1) at (3.2,0)
{\scalebox{0.35}{
\begin{tikzpicture}
%perfect matching
\draw[line width=\pmwidth, color=\pmcolor] (W3)--(B1);
\draw[line width=\pmwidth, color=\pmcolor] (B3)--(6,3); \draw[line width=\pmwidth, color=\pmcolor] (W2)--(0,3);
\draw[line width=\pmwidth, color=\pmcolor] (W1)--(6,0); \draw[line width=\pmwidth, color=\pmcolor] (B2)--(0,6);
\basicdimerB
\end{tikzpicture} }};

\node (PM2) at (6.4,0)
{\scalebox{0.35}{
\begin{tikzpicture}
%perfect matching
\draw[line width=\pmwidth, color=\pmcolor] (W2)--(B1); \draw[line width=\pmwidth, color=\pmcolor] (W3)--(B3);
\draw[line width=\pmwidth, color=\pmcolor] (W1)--(6,0); \draw[line width=\pmwidth, color=\pmcolor] (B2)--(0,6);
\basicdimerB
\end{tikzpicture} }};

\node (PM3) at (9.6,0)
{\scalebox{0.35}{
\begin{tikzpicture}
%perfect matching
\draw[line width=\pmwidth, color=\pmcolor] (W2)--(B2); \draw[line width=\pmwidth, color=\pmcolor] (W1)--(B3);
\draw[line width=\pmwidth, color=\pmcolor] (W3)--(3,6); \draw[line width=\pmwidth, color=\pmcolor] (B1)--(3,0);
\basicdimerB
\end{tikzpicture} }};

\node (PM4) at (12.8,0)
{\scalebox{0.35}{
\begin{tikzpicture}
%perfect matching
\draw[line width=\pmwidth, color=\pmcolor] (W1)--(B1); \draw[line width=\pmwidth, color=\pmcolor] (W3)--(B2);
\draw[line width=\pmwidth, color=\pmcolor] (B3)--(6,3); \draw[line width=\pmwidth, color=\pmcolor] (W2)--(0,3);
\basicdimerB
\end{tikzpicture} }};

\end{tikzpicture}
\caption{Some perfect matchings on the dimer model given in Figure~\ref{ex_dimer_4b}.}\label{pm_4b}
\end{figure}

We say that a dimer model is \emph{non-degenerate} if every edge is contained in some perfect matchings.
It is known that this non-degeneracy condition is equivalent to the \emph{strong marriage condition}; that is, the dimer model has equal numbers of black and white nodes and every proper subset of the black nodes of size $n$ is connected to at least $n+1$ white nodes (e.g., \cite[Remark~2.12]{Bro}).

Following \cite[Section~5]{IU2}, we next define the perfect matching polygon.
We first fix a perfect matching $\sfP_0$, and call this the \emph{reference perfect matching}.
For any perfect matching $\sfP$, we consider the connected components of the universal cover $\RR^2$ divided by $\sfP\cup\sfP_0$.
Then, we consider the \emph{height function} $\sfh_{\sfP,\sfP_0}$, which is a locally constant function on $\RR^2{\setminus}(\sfP\cup\sfP_0)$, defined as follows.
First, we choose a connected component of $\RR^2{\setminus}(\sfP\cup\sfP_0)$, and define the value of $\sfh_{\sfP,\sfP_0}$ as $0$.
Then, this function increases by $1$ when we cross
\begin{itemize}\itemsep=0pt
\item[--] an edge $e\in\sfP$ with the black node on the right, or
\item[--] an edge $e\in\sfP_0$ with the white node on the right,
\end{itemize}
and decreases by $1$
when we cross
\begin{itemize}\itemsep=0pt
\item[--] an edge $e\in\sfP$ with the white node on the right, or
\item[--] an edge $e\in\sfP_0$ with the black node on the right.
\end{itemize}
This function is determined up to a choice of a connected component of value $0$.
For example, Figure~\ref{fig_heightchnge} shows the height function $\sfh_{\sfP_2,\sfP_3}$ on the dimer model given in Figure~\ref{ex_dimer_4b},
where the red square stands for a fundamental domain of $\TT$, the edges in $\sfP_2$ (resp.~$\sfP_3$) are colored blue (resp.\ green),
and the number filled in each component is the value of $\sfh_{\sfP_2,\sfP_3}$.

\begin{figure}[h!]\centering
\newcommand{\edgewidth}{0.07cm}
\newcommand{\nodewidth}{0.07cm}
\newcommand{\noderad}{0.24} %radius_of_node
\newcommand{\arrowwidth}{0.08cm}
\newcommand{\pmwidth}{0.5cm}

\scalebox{0.33}{
\begin{tikzpicture}
%coordinate setting
\coordinate (W1a) at (5,1); \coordinate (W2a) at (1,3); \coordinate (W3a) at (3,5);
\coordinate (B1a) at (3,1); \coordinate (B2a) at (1,5); \coordinate (B3a) at (5,3);
\coordinate (W1b) at (5,7); \coordinate (W2b) at (1,9); \coordinate (W3b) at (3,11);
\coordinate (B1b) at (3,7); \coordinate (B2b) at (1,11); \coordinate (B3b) at (5,9);
\coordinate (W1c) at (11,1); \coordinate (W2c) at (7,3); \coordinate (W3c) at (9,5);
\coordinate (B1c) at (9,1); \coordinate (B2c) at (7,5); \coordinate (B3c) at (11,3);
\coordinate (W1d) at (11,7); \coordinate (W2d) at (7,9); \coordinate (W3d) at (9,11);
\coordinate (B1d) at (9,7); \coordinate (B2d) at (7,11); \coordinate (B3d) at (11,9);

%perfect matching P2
\draw[line width=\pmwidth, color=blue] (W2a)--(B1a); \draw[line width=\pmwidth, color=blue] (W3a)--(B3a);
\draw[line width=\pmwidth, color=blue] (W1a)--(6,0); \draw[line width=\pmwidth, color=blue] (B2a)--(0,6);
\draw[line width=\pmwidth, color=blue] (W2b)--(B1b); \draw[line width=\pmwidth, color=blue] (W3b)--(B3b);
\draw[line width=\pmwidth, color=blue] (W1b)--(B2c); \draw[line width=\pmwidth, color=blue] (B2b)--(0,12);
\draw[line width=\pmwidth, color=blue] (W2c)--(B1c); \draw[line width=\pmwidth, color=blue] (W3c)--(B3c);
\draw[line width=\pmwidth, color=blue] (W1c)--(12,0); \draw[line width=\pmwidth, color=blue] (B2c)--(6,6);
\draw[line width=\pmwidth, color=blue] (W2d)--(B1d); \draw[line width=\pmwidth, color=blue] (W3d)--(B3d);
\draw[line width=\pmwidth, color=blue] (W1d)--(12,6); \draw[line width=\pmwidth, color=blue] (B2d)--(6,12);

%perfect matching P3
\draw[line width=\pmwidth, color=green] (W2a)--(B2a); \draw[line width=\pmwidth, color=green] (W1a)--(B3a);
\draw[line width=\pmwidth, color=green] (W3a)--(B1b); \draw[line width=\pmwidth, color=green] (B1a)--(3,0);
\draw[line width=\pmwidth, color=green] (W2b)--(B2b); \draw[line width=\pmwidth, color=green] (W1b)--(B3b);
\draw[line width=\pmwidth, color=green] (W3b)--(3,12);
\draw[line width=\pmwidth, color=green] (W2c)--(B2c); \draw[line width=\pmwidth, color=green] (W1c)--(B3c);
\draw[line width=\pmwidth, color=green] (W3c)--(B1d); \draw[line width=\pmwidth, color=green] (B1c)--(9,0);
\draw[line width=\pmwidth, color=green] (W2d)--(B2d); \draw[line width=\pmwidth, color=green] (W1d)--(B3d);
\draw[line width=\pmwidth, color=green] (W3d)--(9,12);

%edge
\draw[line width=\edgewidth] (W2a)--(B1a)--(W1a)--(B3a)--(W3a)--(B2a)--(W2a); \draw[line width=\edgewidth] (B1a)--(W3a);
\draw[line width=\edgewidth] (0,3)--(W2a); \draw[line width=\edgewidth] (3,0)--(B1a); \draw[line width=\edgewidth] (6,0)--(W1a);
\draw[line width=\edgewidth] (0,6)--(B2a);
\draw[line width=\edgewidth] (W2b)--(B1b)--(W1b)--(B3b)--(W3b)--(B2b)--(W2b); \draw[line width=\edgewidth] (B1b)--(W3b);
\draw[line width=\edgewidth] (0,9)--(W2b); \draw[line width=\edgewidth] (0,12)--(B2b); \draw[line width=\edgewidth] (3,12)--(W3b);
\draw[line width=\edgewidth] (W2c)--(B1c)--(W1c)--(B3c)--(W3c)--(B2c)--(W2c); \draw[line width=\edgewidth] (B1c)--(W3c);
\draw[line width=\edgewidth] (9,0)--(B1c); \draw[line width=\edgewidth] (12,0)--(W1c); \draw[line width=\edgewidth] (12,3)--(B3c);
\draw[line width=\edgewidth] (W2d)--(B1d)--(W1d)--(B3d)--(W3d)--(B2d)--(W2d); \draw[line width=\edgewidth] (B1d)--(W3d);
\draw[line width=\edgewidth] (12,6)--(W1d);\draw[line width=\edgewidth] (6,12)--(B2d);
\draw[line width=\edgewidth] (12,9)--(B3d); \draw[line width=\edgewidth] (9,12)--(W3d);
\draw[line width=\edgewidth] (W3a)--(B1b); \draw[line width=\edgewidth] (B3a)--(W2c);
\draw[line width=\edgewidth] (B3b)--(W2d); \draw[line width=\edgewidth] (B2c)--(W1b);
\draw[line width=\edgewidth] (W3c)--(B1d);

%black
\draw [line width=\nodewidth, fill=black] (B1a) circle [radius=\noderad] ;
\draw [line width=\nodewidth, fill=black] (B2a) circle [radius=\noderad] ;
\draw [line width=\nodewidth, fill=black] (B3a) circle [radius=\noderad] ;
\draw [line width=\nodewidth, fill=black] (B1b) circle [radius=\noderad] ;
\draw [line width=\nodewidth, fill=black] (B2b) circle [radius=\noderad] ;
\draw [line width=\nodewidth, fill=black] (B3b) circle [radius=\noderad] ;
\draw [line width=\nodewidth, fill=black] (B1c) circle [radius=\noderad] ;
\draw [line width=\nodewidth, fill=black] (B2c) circle [radius=\noderad] ;
\draw [line width=\nodewidth, fill=black] (B3c) circle [radius=\noderad] ;
\draw [line width=\nodewidth, fill=black] (B1d) circle [radius=\noderad] ;
\draw [line width=\nodewidth, fill=black] (B2d) circle [radius=\noderad] ;
\draw [line width=\nodewidth, fill=black] (B3d) circle [radius=\noderad] ;
%white
\draw [line width=\nodewidth, fill=white] (W1a) circle [radius=\noderad] ;
\draw [line width=\nodewidth, fill=white] (W2a) circle [radius=\noderad] ;
\draw [line width=\nodewidth, fill=white] (W3a) circle [radius=\noderad] ;
\draw [line width=\nodewidth, fill=white] (W1b) circle [radius=\noderad] ;
\draw [line width=\nodewidth, fill=white] (W2b) circle [radius=\noderad] ;
\draw [line width=\nodewidth, fill=white] (W3b) circle [radius=\noderad] ;
\draw [line width=\nodewidth, fill=white] (W1c) circle [radius=\noderad] ;
\draw [line width=\nodewidth, fill=white] (W2c) circle [radius=\noderad] ;
\draw [line width=\nodewidth, fill=white] (W3c) circle [radius=\noderad] ;
\draw [line width=\nodewidth, fill=white] (W1d) circle [radius=\noderad] ;
\draw [line width=\nodewidth, fill=white] (W2d) circle [radius=\noderad] ;
\draw [line width=\nodewidth, fill=white] (W3d) circle [radius=\noderad] ;

%height
\node at (1,1) {\Huge$\mathchar`-1$}; \node at (0.35,4) {\Huge$\mathchar`-1$};
\node at (2,3.5) {\Huge$0$}; \node at (4,2.5) {\Huge$0$}; \node at (4,0.5) {\Huge$0$}; \node at (1,7) {\Huge$0$}; \node at (0.35,10) {\Huge$0$};
\node at (5,5) {\Huge$1$}; \node at (7,1) {\Huge$1$}; \node at (4,8.5) {\Huge$1$}; \node at (2,9.5) {\Huge$1$}; \node at (2,11.65) {\Huge$1$};
\node at (10,0.5) {\Huge$2$}; \node at (8,3.5) {\Huge$2$}; \node at (10,2.5) {\Huge$2$}; \node at (7,7) {\Huge$2$}; \node at (5,11) {\Huge$2$};
\node at (11.65,2) {\Huge$3$}; \node at (11,5) {\Huge$3$}; \node at (8,9.5) {\Huge$3$}; \node at (10,8.5) {\Huge$3$}; \node at (8,11.65) {\Huge$3$};
\node at (11,11) {\Huge$4$}; \node at (11.65,8) {\Huge$4$};

%fundamental domain
\draw[line width=0.12cm, red] (0,0) rectangle (6,6);
\end{tikzpicture}
}
\caption{The height function $\sfh_{\sfP_2,\sfP_3}$.}\label{fig_heightchnge}
\end{figure}

We then take a point $\pt\in\RR^2{\setminus}(\sfP\cup\sfP_0)$, and define the \emph{height change}
\begin{displaymath}
h(\sfP,\sfP_0)=(h_x(\sfP,\sfP_0),h_y(\sfP,\sfP_0))\in\ZZ^2
\end{displaymath}
of $\sfP$ with respect to $\sfP_0$ as the differences of the height function:
\begin{gather*}
h_x(\sfP,\sfP_0)=\sfh_{\sfP,\sfP_0}(\pt+(1,0))-\sfh_{\sfP,\sfP_0}(\pt),
\\
h_y(\sfP,\sfP_0)=\sfh_{\sfP,\sfP_0}(\pt+(0,1))-\sfh_{\sfP,\sfP_0}(\pt).
\end{gather*}
We remark that this does not depend on the choice of $\pt\in\RR^2{\setminus}(\sfP\cup\sfP_0)$.
We then consider the height change $h(\sfP,\sfP^\prime)$ for any pair of perfect matchings $\sfP$, $\sfP^\prime$,
but since we have
\begin{equation*}%\label{relation_height}
h(\sfP,\sfP^\prime)=h(\sfP,\sfP_0)-h(\sfP^\prime,\sfP_0),
\end{equation*}
it is enough to consider height changes with respect to the reference perfect matching $\sfP_0$.
Then, the \emph{perfect matching} (= \emph{PM}) \emph{polygon} (or \emph{characteristic polygon}) $\Delta_\Gamma\subset\RR^2$ of a dimer model $\Gamma$ is defined
as the convex hull of $\big\{h(\sfP,\sfP_0)\in\ZZ^2\,|\, \sfP\in\PM(\Gamma)\big\}$ where $\PM(\Gamma)$ is the set of perfect matchings on~$\Gamma$.

\begin{Remark}\label{rem_polygon}
The description of height changes depends on the choice of the coordinate system fixed in $\TT$
(i.e., the choice of a fundamental domain).
A change of a coordinate system induces a $\GL(2,\ZZ)$-action on the PM polygon, and this action does not affect our problem.
In the following, we say that two polygons $P$ and $Q$ are \emph{$\GL(2,\ZZ)$-equivalent} if they are transformed into each other by $\GL(2,\ZZ)$-transformations, in which case we denote $P\cong Q$.
Thus, we may fix some fundamental domain of~$\TT$.
Also, we remark that the description of the polygon $\Delta_\Gamma$ depends on the choice of a reference perfect matching, but it is determined up to translations.
\end{Remark}

\begin{Definition}\label{def_corner}
Fix a perfect matching $\sfP_0$ and let $\Delta_\Gamma$ be the perfect matching polygon.
We say that a perfect matching $\sfP$ is
\begin{itemize}\itemsep=0pt
\item a \emph{corner} (or \emph{extremal}) \emph{perfect matching} if $h(\sfP,\sfP_0)$ is a vertex of $\Delta_\Gamma$,
\item a \emph{boundary} (or \emph{external}) \emph{perfect matching} if $h(\sfP,\sfP_0)$ is a lattice point on an edge of $\Delta_\Gamma$
(in particular, corner perfect matchings are boundary perfect matchings),
\item an \emph{internal} \emph{perfect matching} if $h(\sfP,\sfP_0)$ is an interior lattice point of $\Delta_\Gamma$.
\end{itemize}
\end{Definition}

In the next subsection, we will introduce consistent dimer models (see Definition~\ref{def_consistent}), which have several nice properties.
If a dimer model is consistent, then there exists a unique corner perfect matching corresponding to each vertex of $\Delta_\Gamma$
(e.g., \cite[Corollary~4.27]{Bro}, \cite[Proposition~9.2]{IU2}).
Thus, we can give a cyclic order to corner perfect matchings along the corresponding vertices of $\Delta_\Gamma$ in the anti-clockwise direction.
We say that two corner perfect matchings are \emph{adjacent} if they are adjacent with respect to the given cyclic order.

\begin{Example}
We consider the dimer model given in Figure~\ref{ex_dimer_4b}, which is consistent as we will see in the next subsection.
Fix the perfect matching $\sfP_0$ shown in Figure~\ref{pm_4b} as the reference one.
Then, we see that the perfect matchings $\sfP_1,\dots,\sfP_4$ correspond to lattice points $(1,0)$, $(1,1)$, $(-1,0)$, $(0,-1)$, respectively.
Since $\sfP_0$ is the reference perfect matching, it corresponds to $(0,0)$.
In addition, this dimer model has three perfect matchings that are not listed in Figure~\ref{pm_4b},
and such perfect matchings also correspond to $(0,0)$.
Thus, the PM polygon takes the form shown in Figure~\ref{polygon_4b}, and hence $\sfP_1,\dots,\sfP_4$ are corner perfect matchings.

\begin{figure}[h!]\centering
\scalebox{0.75}{
\begin{tikzpicture}
\coordinate (00) at (0,0); \coordinate (10) at (1,0); \coordinate (01) at (0,1); \coordinate (-10) at (-1,0); \coordinate (11) at (1,1);
\coordinate (0-1) at (0,-1);

\draw [step=1, gray] (-2.3,-2.3) grid (2.3,2.3);
\draw [line width=0.06cm] (10)--(11) ; \draw [line width=0.06cm] (11)--(-10) ;
\draw [line width=0.06cm] (-10)--(0-1) ; \draw [line width=0.06cm] (0-1)--(10) ;

\draw [line width=0.05cm, fill=black] (00) circle [radius=0.1] ; \draw [line width=0.05cm, fill=black] (10) circle [radius=0.1] ;
\draw [line width=0.05cm, fill=black] (11) circle [radius=0.1] ;
\draw [line width=0.05cm, fill=black] (-10) circle [radius=0.1] ; \draw [line width=0.05cm, fill=black] (0-1) circle [radius=0.1] ;

\node [red] at (1.5,0) {$\sfP_1$} ;
\node [red] at (1.5,1) {$\sfP_2$} ; \node [red] at (-1.5,0) {$\sfP_3$} ;
\node [red] at (0,-1.5) {$\sfP_4$} ;

\end{tikzpicture}
}
\caption{The PM polygon of the dimer model given in Figure~\ref{ex_dimer_4b}.}\label{polygon_4b}
\end{figure}
\end{Example}

In this way, we can obtain the PM polygon from a dimer model.
On the other hand, it is known that any lattice polygon can be obtained as the PM polygon of a certain dimer model.

\begin{Theorem}[\cite{Gul,IU2}]\label{existence_dimer}
For any lattice polygon $\Delta$ in $\RR^2$, there exists a dimer model $\Gamma$ giving $\Delta$ as the PM polygon $\Delta_\Gamma$.
Furthermore, we can take this $\Gamma$ as it satisfies the consistency condition $($see Definition~{\rm \ref{def_consistent})}.
\end{Theorem}

Thus, for a given lattice polygon $\Delta$, we say that $\Gamma$ is a \emph{dimer model associated with $\Delta$} if the PM polygon of~$\Gamma$ coincides with~$\Delta$.
We remark that for a given polygon $\Delta$, the associated consistent dimer model is not unique in general.

\section{Zigzag paths and their properties}\label{sec_zigzag}

\subsection{Consistency conditions}\label{subsec_consistency}

In this subsection, we introduce the consistency condition.
In order to define this condition, we first introduce the notion of zigzag paths.
These paths are also the main ingredients for introducing deformations of dimer models.

\begin{Definition}We say that a path on a dimer model is a \emph{zigzag path} if it makes a maximum turn to the right on a white node and a maximum turn to the left on a black node.
Also, we say that a zigzag path is \emph{reduced} if it does not pass through $2$-valent nodes.
(We remark that we can make a zigzag path reduced using the join moves. In particular, any zigzag path on a~reduced dimer model is reduced.)
\end{Definition}

Since a dimer model has only finitely many edges, we see that all zigzag paths are periodic.
For a zigzag path $z$ on $\Gamma$, we define the length of~$z$, which is denoted by~$\ell(z)$, as the number of edges of~$\Gamma$ constituting~$z$.
In particular, we see that $\ell(z)$ is an even integer.
Thus, edges on a~zigzag path are indexed by elements in $\ZZ/(2n)\ZZ$ for some integer $\ell(z)/2=n\ge 1$.
Fix a black node on a zigzag path $z$ as the starting point of~$z$, and denote~$z$ as a sequence of edges starting from the fixed black node:
$z=z[1]z[2]\cdots z[2n-1]z[2n]$.

\begin{center}
{\scalebox{0.7}{
\begin{tikzpicture}
\newcommand{\edgewidth}{0.05cm} %line_width_of _edges
\newcommand{\nodewidth}{0.05cm} %line_width_around_nodes
\newcommand{\noderad}{0.16} %radius_of_node

%coordinate setting
\coordinate (B1) at (0,0); \coordinate (B2) at (3,0); \coordinate (B3) at (6,0);
\coordinate (W1) at (1.5,1.5); \coordinate (W2) at (4.5,1.5); \coordinate (W3) at (7.5,1.5);

\path (B1) ++(135:1.5cm) coordinate (B1w); \path (B1) ++(225:0.7cm) coordinate (B1s); \path (B1) ++(315:0.7cm) coordinate (B1e);
\path (W1) ++(45:0.7cm) coordinate (W1w); \path (W1) ++(135:0.7cm) coordinate (W1s);
\path (B2) ++(225:0.7cm) coordinate (B2s); \path (B2) ++(315:0.7cm) coordinate (B2e);
\path (W2) ++(45:0.7cm) coordinate (W2w); \path (W2) ++(135:0.7cm) coordinate (W2s);
\path (B3) ++(225:0.7cm) coordinate (B3s); \path (B3) ++(315:0.7cm) coordinate (B3e);
\path (W3) ++(315:1.5cm) coordinate (W3n); \path (W3) ++(45:0.7cm) coordinate (W3w); \path (W3) ++(135:0.7cm) coordinate (W3s);

\path (B1) ++(245:0.5cm) coordinate (B1ss); \path (B1) ++(295:0.5cm) coordinate (B1ee);
\path (W1) ++(65:0.5cm) coordinate (W1ww); \path (W1) ++(115:0.5cm) coordinate (W1ss);
\path (B2) ++(245:0.5cm) coordinate (B2ss); \path (B2) ++(295:0.5cm) coordinate (B2ee);
\path (W2) ++(65:0.5cm) coordinate (W2ww); \path (W2) ++(115:0.5cm) coordinate (W2ss);
\path (B3) ++(245:0.5cm) coordinate (B3ss); \path (B3) ++(295:0.5cm) coordinate (B3ee);
\path (W3) ++(65:0.5cm) coordinate (W3ww); \path (W3) ++(115:0.5cm) coordinate (W3ss);
%edge
\draw [line width=\edgewidth] (B1w)--(B1)--(W1)--(B2)--(W2)--(B3)--(W3)--(W3n);
\draw [line width=\edgewidth] (B1s)--(B1)--(B1e); \draw[line width=\edgewidth] (B2s)--(B2)--(B2e); \draw [line width=\edgewidth] (B3s)--(B3)--(B3e);
\draw [line width=\edgewidth] (W1w)--(W1)--(W1s); \draw[line width=\edgewidth] (W2w)--(W2)--(W2s); \draw [line width=\edgewidth] (W3w)--(W3)--(W3s);

\draw [line width=\edgewidth,line cap=round, dash pattern=on 0pt off 2.5\pgflinewidth] (B1ss)--(B1ee);
\draw [line width=\edgewidth,line cap=round, dash pattern=on 0pt off 2.5\pgflinewidth] (W1ww)--(W1ss);
\draw [line width=\edgewidth,line cap=round, dash pattern=on 0pt off 2.5\pgflinewidth] (B2ss)--(B2ee);
\draw [line width=\edgewidth,line cap=round, dash pattern=on 0pt off 2.5\pgflinewidth] (W2ww)--(W2ss);
\draw [line width=\edgewidth,line cap=round, dash pattern=on 0pt off 2.5\pgflinewidth] (B3ss)--(B3ee);
\draw [line width=\edgewidth,line cap=round, dash pattern=on 0pt off 2.5\pgflinewidth] (W3ww)--(W3ss);
%blacknode
\draw [line width=\nodewidth, fill=black] (B1) circle [radius=\noderad] ;
\draw [line width=\nodewidth, fill=black] (B2) circle [radius=\noderad] ;
\draw [line width=\nodewidth, fill=black] (B3) circle [radius=\noderad] ;
%whitenode
\draw [line width=\nodewidth, fill=white] (W1) circle [radius=\noderad] ;
\draw [line width=\nodewidth, fill=white] (W2) circle [radius=\noderad] ;
\draw [line width=\nodewidth, fill=white] (W3) circle [radius=\noderad] ;
%zigzag
\path (B1w) ++(90:0.2cm) coordinate (B1w+);
\path (B1) ++(90:0.2cm) coordinate (B1+); \path (W1) ++(90:0.2cm) coordinate (W1+);
\path (B2) ++(90:0.2cm) coordinate (B2+); \path (W2) ++(90:0.2cm) coordinate (W2+);
\path (B3) ++(90:0.2cm) coordinate (B3+); \path (W3) ++(90:0.2cm) coordinate (W3+);
\path (W3n) ++(90:0.2cm) coordinate (W3n+);
\draw [->, rounded corners, line width=0.08cm, red] (B1w+)--(B1+)--(W1+)--(B2+)--(W2+)--(B3+)--(W3+)--(W3n+) ;
\draw [line width=0.05cm, red, line cap=round, dash pattern=on 0pt off 2.5\pgflinewidth] (-1.8,0.8)--(-2.4,0.8) ;
\draw [line width=0.05cm, red, line cap=round, dash pattern=on 0pt off 2.5\pgflinewidth] (9.3,0.8)--(9.9,0.8) ;
%index
\node[red] at (0.5,1.3) {\large$z[1]$} ; \node[red] at (2,0.5) {\large$z[2]$} ;
\node[red] at (3.5,1.3) {\large$z[3]$} ; \node[red] at (5,0.5) {\large$z[4]$} ;
\node[red] at (6.5,1.3) {\large$z[5]$} ; \node[red] at (8,0.5) {\large$z[6]$} ;
\node[red] at (-1.1,0.5) {\large$z[2n]$} ;
\end{tikzpicture}
} }
\end{center}

An edge in a zigzag path $z$ is called a \emph{zig} (resp.\ \emph{zag}) of $z$ if it is indexed by an odd (resp.~even) integer.
We denote by $\Zig(z)$ (resp.~$\Zag(z)$) the set of zigs (resp.~zags) appearing in a~zigzag path~$z$, which is a~finite set.
Two zigzag paths are said to \emph{intersect} if they share an edge (not a node).
We note that if $z$ does not have a self-intersection, $\Zig(z)$ and $\Zag(z)$ are disjoint sets.
For any edge~$e$ of a dimer model, we can consider the zigzag path containing $e$ as a zig and the zigzag path containing $e$ as a zag.
Thus, any edge $e$ is contained in at most two zigzag paths.
If such zigzag paths do not have a self-intersection, $e$ is contained in exactly two zigzag paths.
For example, zigzag paths on the dimer model given in the left of Figure~\ref{ex_dimer_4b} are shown in Figure~\ref{zigzag_4b}.

\begin{figure}[h!]
\centering

\newcommand{\edgewidth}{0.07cm}
\newcommand{\nodewidth}{0.07cm}
\newcommand{\noderad}{0.24} %radius_of_node

{\scalebox{1}{
\begin{tikzpicture}

\node at (0,-1.6) {$z_1$}; \node at (3.5,-1.6) {$z_2$}; \node at (7,-1.6) {$z_3$}; \node at (10.5,-1.6) {$z_4$};

\node (ZZ1) at (0,0)
{\scalebox{0.4}{
\begin{tikzpicture}
\basicdimerB
%zigzag path
\draw[->, line width=0.2cm, rounded corners, color=red] (0,3)--(W2)--(B1)--(W3)--(B3)--(6,3);
\end{tikzpicture} }};

\node (ZZ2) at (3.5,0)
{\scalebox{0.4}{
\begin{tikzpicture}
\basicdimerB
%zigzag path
\draw[->, line width=0.2cm, rounded corners, color=red] (3,0)--(B1)--(W2)--(B2)--(0,6);
\draw[->, line width=0.2cm, rounded corners, color=red] (6,0)--(W1)--(B3)--(W3)--(3,6);
\end{tikzpicture} }} ;

\node (ZZ3) at (7,0)
{\scalebox{0.4}{
\begin{tikzpicture}
\basicdimerB
%zigzag path
\draw[->, line width=0.2cm, rounded corners, color=red] (3,6)--(W3)--(B2)--(W2)--(0,3);
\draw[->, line width=0.2cm, rounded corners, color=red] (6,3)--(B3)--(W1)--(B1)--(3,0);
\end{tikzpicture} }};

\node (ZZ4) at (10.5,0)
{\scalebox{0.4}{
\begin{tikzpicture}
\basicdimerB
%zigzag path
\draw[->, line width=0.2cm, rounded corners, color=red] (0,6)--(B2)--(W3)--(B1)--(W1)--(6,0);
\end{tikzpicture} }} ;

\end{tikzpicture}
}}
\caption{All zigzag paths on the dimer model given in Figure~\ref{ex_dimer_4b}.}\label{zigzag_4b}
\end{figure}

For a zigzag path $z$ on a dimer model $\Gamma$, we also consider the lift of $z$ to the universal cover~$\widetilde{\Gamma}$.
Let $\widetilde{z}{(\alpha)}$ denote a zigzag path on $\widetilde{\Gamma}$ whose projection on $\Gamma$ is $z$ where $\alpha\in\ZZ$.
When we do not need to specify these, we simply denote each of them by~$\widetilde{z}$.
Then, we see that a zigzag path on~$\widetilde{\Gamma}$ is either periodic or infinite in both directions.
Using these notions, we introduce the consistency condition.

\begin{Definition}[{see \cite[Definition~3.5]{IU1}}]
\label{def_consistent}
We say that a dimer model is (\emph{zigzag}) \emph{consistent} if it satisfies the following conditions:
\begin{itemize}\itemsep=0pt
\item [(1)] there is no homologically trivial zigzag path,
\item [(2)] no zigzag path on the universal cover has a self-intersection,
\item [(3)] no pair of zigzag paths on the universal cover intersect each other in the same direction more than once.
That is, if a pair of zigzag paths $(\widetilde{z},\widetilde{w})$ on the universal cover has two intersected edges $a_1$, $a_2$
and $\widetilde{z}$ points from $a_1$ to $a_2$, then $\widetilde{w}$ points from $a_2$ to $a_1$.
\end{itemize}
\end{Definition}

In the literature, there are several conditions that are equivalent to Definition~\ref{def_consistent} (for more details, see \cite{Boc1,IU1}),
and it is known that a consistent dimer model is non-degenerate (e.g., \cite[Proposition~8.1]{IU2}).
For example, we see that the dimer model given in Figure~\ref{ex_dimer_4b} is consistent by checking all zigzag paths, which are shown in Figure~\ref{zigzag_4b}.
We also remark that this dimer model satisfies the stronger condition called \emph{isoradial} (see Definition~\ref{def_isoradial}).

In this paper, we also use another condition called \emph{properly ordered}.
To explain the proper ordering, we prepare several notation.
First, considering a zigzag path $z$ as a $1$-cycle on $\TT$, we have the homology class $[z]\in\rmH_1(\TT)\cong\ZZ^2$. We call this element $[z]\in\ZZ^2$ the \emph{slope} of $z$.
We remark that even if we apply the join and split moves to nodes contained in a zigzag path, such operations do not change the slope.
If a zigzag path does not have any self-intersection, the slope of each zigzag path is primitive.
Now, we consider slopes $(a,b)\in\ZZ^2$ of zigzag paths that are not homologically trivial.
The set of such slopes has a natural cyclic order by considering~$(a,b)$ as an element of the unit circle:
\begin{displaymath}
\frac{(a,b)}{\sqrt{a^2+b^2}}\in S^1.
\end{displaymath}
We say that two zigzag paths are \emph{adjacent} if their slopes are adjacent with respect to the above cyclic order.
Using this cyclic order, we define a properly ordered dimer model below.
In particular, it is known that a dimer model is consistent in the sense of Definition~\ref{def_consistent} if and only if it is properly ordered (see \cite[Proposition~4.4]{IU1}).

\begin{Definition}[{see \cite[Section~3.1]{Gul}}]\label{def_properly}
A dimer model is said to be \emph{properly ordered} if
\begin{itemize}\itemsep=0pt
\item [(1)] there is no homologically trivial zigzag path,
\item [(2)] no zigzag path on the universal cover has a self-intersection,
\item [(3)] no pair of zigzag paths with the same slope have a common node,
\item [(4)] for any node on the dimer model, the natural cyclic order on the set of zigzag paths incident to that node
coincides with the cyclic order determined by their slopes.
\end{itemize}
\end{Definition}

We also introduce isoradial dimer models which are stronger than consistent ones.
The dimer model given in Figure~\ref{ex_dimer_4b} is isoradial in particular.

\begin{Definition}[{\cite[Theorem~5.1]{KS}; see also \cite{Duf,Mer}}]
\label{def_isoradial}
We say that a dimer model $\Gamma$ is \emph{isoradial} (or \emph{geometrically consistent}) if
\begin{itemize}\itemsep=0pt
\item [(1)] every zigzag path is a simple closed curve,
\item [(2)] any pair of zigzag paths on the universal cover share at most one edge.
\end{itemize}
\end{Definition}

\subsection{Relationships between perfect matchings and zigzag paths}\label{subsec_relation_zigzag_pm}

We can now discuss the relationship between perfect matchings and zigzag paths.
The following proposition is essential throughout this paper.

\begin{Proposition}[{see \cite[Theorem~3.3 and Corollary~3.8]{Gul}, \cite[Proposition~9.2 and Corollary~9.3]{IU2}}]\label{zigzag_sidepolygon}
There exists a one-to-one correspondence between the set of slopes of zigzag paths on a consistent dimer model $\Gamma$ and
the set of primitive side segments of the PM polygon $\Delta_\Gamma$. More precisely, each slope of a zigzag path is the primitive outer normal vector for a primitive side segment of~$\Delta_\Gamma$.

Moreover, zigzag paths having the same slope arise as the difference of two adjacent corner perfect matchings $\sfP$, $\sfP^\prime$ $($i.e., the edges in $\sfP\cup\sfP^\prime{\setminus}\sfP\cap\sfP^\prime$ form zigzag paths$)$. Thus, any corner perfect matching intersects with half of the edges constituting a certain zigzag path.
\end{Proposition}

For example, the zigzag path $z_1$ shown in Figure~\ref{zigzag_4b} is obtained from the pair of adjacent corner perfect matchings $(\sfP_1,\sfP_2)$ given in Figure~\ref{pm_4b}. Also, the zigzag paths $z_2$, $z_3$, $z_4$ are obtained by pairs $(\sfP_2,\sfP_3)$, $(\sfP_3,\sfP_4)$, $(\sfP_4,\sfP_1)$, respectively.

By Proposition~\ref{zigzag_sidepolygon}, we can assign each edge of the PM polygon to a zigzag path~$z$;
thus we will call this the edge corresponding to~$z$.
In particular, the edges corresponding to zigzag paths having the same slope are all the same.

Let $\sfP$, $\sfP^\prime$ be adjacent corner perfect matchings on a consistent dimer model, and $z_1,\dots, z_r$ be the zigzag paths arising from $\sfP$ and $\sfP^\prime$ as in Proposition~\ref{zigzag_sidepolygon}.
In particular, these zigzag paths have the same slope.
We see that $\sfP\cap z_i=\Zig(z_i)$ and $\sfP^\prime\cap z_i=\Zag(z_i)$ (or $\sfP\cap z_i=\Zag(z_i)$ and $\sfP^\prime\cap z_i=\Zig(z_i)$) for any $i=1,\dots,r$. Here, $\sfP\cap z_i$ denotes the subset of edges in $\sfP$ contained in~$z_i$.
Then, we have the description of boundary perfect matchings using the corner ones.

\begin{Proposition}[{e.g., \cite[Proposition~4.35]{Bro}, \cite[Corollary~3.8]{Gul}}]\label{char_bound}
Let $\sfP$, $\sfP^\prime$ and $z_1,\dots, z_r$ be as above.
Let $E$ be the edge of the PM polygon of $\Gamma$ corresponding to $z_1,\dots, z_r$.
We assume that $\sfP\cap z_i=\Zig(z_i)$ and $\sfP^\prime\cap z_i=\Zag(z_i)$.
Then, any boundary perfect matching corresponding to a~lattice point on $E$ can be described as
\begin{displaymath}
\bigg(\sfP{\setminus}\bigcup_{i\in I}\Zig(z_i)\bigg)\cup\bigcup_{i\in I}\Zag(z_i) \qquad or \qquad \bigg(\sfP^\prime{\setminus}\bigcup_{i\in I}\Zag(z_i)\bigg)\cup\bigcup_{i\in I}\Zig(z_i),
\end{displaymath}
where $I$ is a subset of $\{1, \dots, r\}$.
In particular, the number of perfect matchings corresponding to a~lattice point~$\mathsf{q}$ on~$E$ is~$\dbinom{r}{m}$,
where~$m$ is the number of primitive side segments of~$E$ between~$\mathsf{q}$ and one of the endpoints of~$E$.
\end{Proposition}

We then observe the relationships between zigzag paths and height changes of perfect matchings.
Some of these relationships are well-known to experts, but we will consider them in detail because these statements are quite important
when we define the deformation of consistent dimer models in Section~\ref{sec_def_deform}, and also for the self-containedness.

\begin{observation}[{cf.~\cite[Section~5.3]{IU2}}]\label{height_zigzag}
Let $\Gamma$ be a consistent dimer model.
For any zigzag path~$z$, the slope $[z]$ is an element in $\rmH_1(\TT)$.
On the other hand, we can consider height changes as elements in the cohomology group $\rmH^1(\TT)\cong\ZZ^2$.
We consider a pairing \mbox{$\langle{-},{-}\rangle\colon \! \rmH^1(\TT)\!\times\!\rmH_1(\TT)\!\rightarrow\!\ZZ$}.
By Propositions~\ref{zigzag_sidepolygon} and \ref{char_bound}, there is a perfect matching $\sfP^\prime$ that intersects half of the edges constituting~$z$.
Then, for any perfect matching $\sfP$, we have $\langle h(\sfP,\sfP^\prime),[z]\rangle\le 0$.
In fact, we first replace~$z$ by the path~$p_z$ on the quiver $Q_\Gamma$ going along the left side of~$z$ (see the figure below).

\begin{center}
{\scalebox{0.6}{
\begin{tikzpicture}[sarrow/.style={black, -latex, very thick},tarrow/.style={black, latex-, very thick}]
\newcommand{\edgewidth}{0.05cm} %line_width_of _edges
\newcommand{\nodewidth}{0.05cm} %line_width_around_nodes
\newcommand{\noderad}{0.16} %radius_of_node

%coordinate setting
\coordinate (B1) at (0,0); \coordinate (B2) at (3,0); \coordinate (B3) at (6,0);
\coordinate (W1) at (1.5,1.5); \coordinate (W2) at (4.5,1.5); \coordinate (W3) at (7.5,1.5);

\path (B1) ++(135:1.5cm) coordinate (B1w); \path (B1) ++(225:0.7cm) coordinate (B1s); \path (B1) ++(315:0.7cm) coordinate (B1e);
\path (W1) ++(45:1.5cm) coordinate (W1w); \path (W1) ++(135:1.5cm) coordinate (W1s);
\path (B2) ++(225:0.7cm) coordinate (B2s); \path (B2) ++(315:0.7cm) coordinate (B2e);
\path (W2) ++(45:1.5cm) coordinate (W2w); \path (W2) ++(135:1.5cm) coordinate (W2s);
\path (B3) ++(225:0.7cm) coordinate (B3s); \path (B3) ++(315:0.7cm) coordinate (B3e);
\path (W3) ++(315:1.5cm) coordinate (W3n); \path (W3) ++(45:1.5cm) coordinate (W3w); \path (W3) ++(135:1.5cm) coordinate (W3s);

\path (B1) ++(245:0.5cm) coordinate (B1ss); \path (B1) ++(295:0.5cm) coordinate (B1ee);
\path (W1) ++(65:0.5cm) coordinate (W1ww); \path (W1) ++(115:0.5cm) coordinate (W1ss);
\path (B2) ++(245:0.5cm) coordinate (B2ss); \path (B2) ++(295:0.5cm) coordinate (B2ee);
\path (W2) ++(65:0.5cm) coordinate (W2ww); \path (W2) ++(115:0.5cm) coordinate (W2ss);
\path (B3) ++(245:0.5cm) coordinate (B3ss); \path (B3) ++(295:0.5cm) coordinate (B3ee);
\path (W3) ++(65:0.5cm) coordinate (W3ww); \path (W3) ++(115:0.5cm) coordinate (W3ss);
%edge
\draw [line width=\edgewidth] (B1w)--(B1)--(W1)--(B2)--(W2)--(B3)--(W3)--(W3n);
\draw [line width=\edgewidth] (B1s)--(B1)--(B1e); \draw[line width=\edgewidth] (B2s)--(B2)--(B2e); \draw [line width=\edgewidth] (B3s)--(B3)--(B3e);
\draw [line width=\edgewidth] (W1w)--(W1)--(W1s); \draw[line width=\edgewidth] (W2w)--(W2)--(W2s); \draw [line width=\edgewidth] (W3w)--(W3)--(W3s);

\draw [line width=\edgewidth,line cap=round, dash pattern=on 0pt off 2.5\pgflinewidth] (B1ss)--(B1ee);
\draw [line width=\edgewidth,line cap=round, dash pattern=on 0pt off 2.5\pgflinewidth] (W1ww)--(W1ss);
\draw [line width=\edgewidth,line cap=round, dash pattern=on 0pt off 2.5\pgflinewidth] (B2ss)--(B2ee);
\draw [line width=\edgewidth,line cap=round, dash pattern=on 0pt off 2.5\pgflinewidth] (W2ww)--(W2ss);
\draw [line width=\edgewidth,line cap=round, dash pattern=on 0pt off 2.5\pgflinewidth] (B3ss)--(B3ee);
\draw [line width=\edgewidth,line cap=round, dash pattern=on 0pt off 2.5\pgflinewidth] (W3ww)--(W3ss);
%blacknode
\draw [line width=\nodewidth, fill=black] (B1) circle [radius=\noderad] ;
\draw [line width=\nodewidth, fill=black] (B2) circle [radius=\noderad] ;
\draw [line width=\nodewidth, fill=black] (B3) circle [radius=\noderad] ;
%whitenode
\draw [line width=\nodewidth, fill=white] (W1) circle [radius=\noderad] ;
\draw [line width=\nodewidth, fill=white] (W2) circle [radius=\noderad] ;
\draw [line width=\nodewidth, fill=white] (W3) circle [radius=\noderad] ;
%zigzag
\path (B1w) ++(-90:0.2cm) coordinate (B1w+);
\path (B1) ++(-90:0.2cm) coordinate (B1+); \path (W1) ++(-90:0.2cm) coordinate (W1+);
\path (B2) ++(-90:0.2cm) coordinate (B2+); \path (W2) ++(-90:0.2cm) coordinate (W2+);
\path (B3) ++(-90:0.2cm) coordinate (B3+); \path (W3) ++(-90:0.2cm) coordinate (W3+);
\path (W3n) ++(-90:0.2cm) coordinate (W3n+);
\draw [->, rounded corners, line width=0.08cm, red] (B1w+)--(B1+)--(W1+)--(B2+)--(W2+)--(B3+)--(W3+)--(W3n+) ;

%arrow
\path (B1)++(90:1.5cm) coordinate (V1);
\path (B2)++(90:1.5cm) coordinate (V2);
\path (B3)++(90:1.5cm) coordinate (V3);
\path (9,0)++(90:1.5cm) coordinate (V4);
\draw [sarrow, line width=0.08cm, blue] (V1)-- ++(45:1.7cm);
\draw [sarrow, line width=0.08cm, blue] (V2)-- ++(45:1.7cm);
\draw [tarrow, line width=0.08cm, blue] (V2)-- ++(135:1.7cm);
\draw [sarrow, line width=0.08cm, blue] (V3)-- ++(45:1.7cm);
\draw [tarrow, line width=0.08cm, blue] (V3)-- ++(135:1.7cm);
\draw [tarrow, line width=0.08cm, blue] (V4)-- ++(135:1.7cm);

\path (W1) ++(75:1cm) coordinate (W1www); \path (W1) ++(105:1cm) coordinate (W1sss);
\path (W2) ++(75:1cm) coordinate (W2www); \path (W2) ++(105:1cm) coordinate (W2sss);
\path (W3) ++(75:1cm) coordinate (W3www); \path (W3) ++(105:1cm) coordinate (W3sss);
\draw [line width=\edgewidth,line cap=round, dash pattern=on 0pt off 2.5\pgflinewidth,blue] (W1www)--(W1sss);
\draw [line width=\edgewidth,line cap=round, dash pattern=on 0pt off 2.5\pgflinewidth,blue] (W2www)--(W2sss);
\draw [line width=\edgewidth,line cap=round, dash pattern=on 0pt off 2.5\pgflinewidth,blue] (W3www)--(W3sss);
%node
\node[red] at (9.3,0.35) {\LARGE $z$} ; \node[blue] at (9.3,1.35) {\LARGE $p_z$} ;
\end{tikzpicture}
} }
\end{center}


Then, considering this path $p_z$ as the element $[p_z]\in\rmH_1(\TT)$, we have $[z]=[p_z]$.
By a choice of $\sfP^\prime$, this $p_z$ does not cross any edge in $\sfP^\prime$, and if $p_z$ crosses an edge in $\sfP$, we can see the white node on the right by the definition of $Q_\Gamma$. Thus, we have the desired inequality.
\end{observation}

For a perfect matching $\sfP$ and a zigzag path $z$ on a dimer model $\Gamma$, we denote by $|\sfP\cap z|$ the number of edges in $\sfP\cap z$.
Since the number of perfect matchings is finite, the maximum (resp.\ minimum) number $\omega_{\max}(z)$ (resp.\ $\omega_{\min}(z)$) of $|\sfP\cap z|$ exists for each zigzag path $z$.
For a~consistent dimer model, $z$ can be obtained as the difference of adjacent perfect matchings (see Proposition~\ref{char_bound});
thus we clearly have $\ell(z)/2=\omega_{\max}(z)$.
We set
\begin{gather*}
\PM_{\max}(z)=\{\sfP\in\PM(\Gamma) \,|\, |\sfP\cap z|=\omega_{\max}(z)\},
\\
\PM_{\min}(z)=\{\sfP\in\PM(\Gamma) \,|\, |\sfP\cap z|=\omega_{\min}(z)\}.
\end{gather*}
In particular, if $\sfP$, $\sfP^\prime$ are adjacent corner perfect matchings on a consistent dimer model $\Gamma$, and $z$ is one of the zigzag paths obtained by $\sfP$, $\sfP^\prime$, then we have $\sfP\cap z=\Zig(z)$ and $\sfP^\prime\cap z=\Zag(z)$ (or $\sfP\cap z=\Zag(z)$ and $\sfP^\prime\cap z=\Zig(z)$),
and hence the next lemma easily follows from Propositions~\ref{zigzag_sidepolygon} and~\ref{char_bound}.

\begin{Lemma}\label{number_zigzag_corner}
Let $z$ be a zigzag path on a consistent dimer model $\Gamma$, and $E$ be the edge of the PM polygon of~$\Gamma$ corresponding to~$z$.
If $\sfP_1,\dots,\sfP_s$ are boundary perfect matchings corresponding to lattice points on~$E$, then we have
$\{\sfP_1,\dots,\sfP_s\}=\PM_{\max}(z)$, and $\omega_{\max}(z)=|\sfP_i\cap z|=\ell(z)/2$ for any $i=1,\dots,s$.
\end{Lemma}

Next, we prepare several lemmas, which play crucial roles to define deformations of consistent dimer models.

\begin{Lemma}\label{zigzag_lem1}
Let the notation be the same as Lemma~{\rm \ref{number_zigzag_corner}}.
For any perfect matching~$\sfP$, we have
\begin{displaymath}
|\sfP\cap z|=\ell(z)/2-\langle h(\sfP,\sfP_i),-[z]\rangle.
\end{displaymath}
In particular, we have
\begin{displaymath}
\langle h(\sfP,\sfP_i),-[z]\rangle\le \omega_{\max}(z)-\omega_{\min}(z),
\end{displaymath}
and the equality holds for $\sfP\in \PM_{\min}(z)$.
\end{Lemma}

\begin{proof}
First, the maximum number of $|\sfP\cap z|$ is $\ell(z)/2$, in which case $\sfP=\sfP_i$ by Lemma~\ref{number_zigzag_corner}.
If the path $p_z$ as in Observation~\ref{height_zigzag} crosses an edge $e$ in $\sfP$, it means that $e$ is not an edge constituting $z$, and thus any edge sharing the same white node as $e$ is not contained in $\sfP$.
By Observation~\ref{height_zigzag}, we see that for any perfect matching $\sfP$ the number of edges in $\sfP$ intersecting with $p_z$ coincides with $-\langle h(\sfP,\sfP_i),[z]\rangle$, and thus we have the first equation.

The second assertion follows from the first equation and Lemma~\ref{number_zigzag_corner}.
\end{proof}

By this lemma, we see that $\sfP\in\PM_{\min}(z)$ if and only if $\langle h(\sfP,\sfP_i),[z]\rangle{\le} \langle h(\sfP^\prime,\sfP_i),[z]\rangle$ for any $\sfP^\prime\in\PM(\Gamma)$.
Thus, we see that $\sfP\in\PM_{\min}(z)$ lies either on a vertex of the PM polygon $\Delta_\Gamma$ or an edge of $\Delta_\Gamma$.
Also, even if two zigzag paths $z_j$ and $z_k$ have the same slope, $\ell(z_j)\neq\ell(z_k)$ and $|\sfP\cap z_j|\neq|\sfP\cap z_k|$ in general,
but the difference of these values is the same in the following sense:

\begin{Lemma}\label{zigzag_lem2}
Let $\Gamma$ be a consistent dimer model, and $z$, $z^\prime$ be zigzag paths on $\Gamma$ having the same slope.
Then, for any perfect matching $\sfP$, we have
\begin{displaymath}
\ell(z)/2-|\sfP\cap z|=\ell(z^\prime)/2-|\sfP\cap z^\prime|.
\end{displaymath}
In particular, we have
\begin{displaymath}
\ell(z)/2-\omega_{\min}(z)=\ell(z^\prime)/2-\omega_{\min}(z^\prime).
\end{displaymath}
\end{Lemma}

\begin{proof}Since $[z]=[z^\prime]$, the first equation follows from Lemma~\ref{zigzag_lem1}.
Considering a perfect matching $\sfP$ such that the value of $\langle h(\sfP,\sfP_i),-[z]\rangle=\langle h(\sfP,\sfP_i),-[z^\prime]\rangle$ is maximal, we have the second equation.
\end{proof}

We then divide zigzag paths on a consistent dimer model into the following two types. Note that type I zigzag paths are used to define the deformation of consistent dimer models.

\begin{Definition}\label{def_zigzag_type}
Let $\Gamma$ be a dimer model, and $z$ be a zigzag path on $\Gamma$.
\begin{itemize}\itemsep=0pt
\item[(1)] We say that $z$ is \emph{type I} if $z$ is reduced and $\widetilde{z}$ intersects with any other zigzag path on the universal cover $\widetilde{\Gamma}$ at most once.
\item[(2)] We say that $z$ is \emph{type II} if $z$ is reduced and there exists a zigzag path $\widetilde{w}$ on the universal cover $\widetilde{\Gamma}$ such that $\widetilde{w}$ intersects with $\widetilde{z}$ in the opposite direction more than once.
\end{itemize}
We note that any zigzag path on a reduced consistent dimer model is either type~I or type~II.
In particular, if $\Gamma$ is isoradial, then all zigzag paths are type~I (see Definition~\ref{def_isoradial}).
\end{Definition}

As the following lemmas show, the properties of type I zigzag paths are particularly nice.

\begin{Lemma}\label{lem_existence_pm}
Let $z$ be a type I zigzag path on a consistent dimer model $\Gamma$.
Then, there exists a perfect matching $\sfP$ on $\Gamma$ satisfying $|\sfP\cap z|=0$, which means that $\sfP$ is in $\PM_{\min}(z)$.
\end{Lemma}

\begin{proof}
In order to find a perfect matching $\sfP$, we will use the method discussed in \cite[Section~3]{Gul}, \cite[Section~4]{Bro}.
For this, we first prepare some notation.

We consider the sequence $[z_1],\dots,[z_n]$ of slopes of zigzag paths on~$\Gamma$.
Since $\Gamma$ is consistent, it is properly ordered, and thus we can assume that the slopes are ordered cyclically with this order.
We note that some of the slopes may coincide.
Then, we define the normal fan in $\rmH_1(\TT)\otimes_\ZZ\RR$ whose rays are slopes $[z_1],\dots,[z_n]$.
In particular, each two-dimensional cone $\sigma$ is generated by different adjacent slopes.
We denote by $\theta_i$ the angle formed by $[z_i]$. Here, we suppose that $z=z_k$.
Let $\calR$ be a ray whose angle is $\theta_k+\pi+\epsilon$ where $\epsilon>0$ is a sufficiently small angle satisfying the condition that $\theta_k+\pi+\epsilon$ does not coincide with any $\theta_i$.

Then, for each node $v\in\Gamma_0$, we define the fan $\xi(v)$ generated by the slopes of zigzag paths factoring through $v$.
In this fan $\xi(v)$, we can find the zigzag path $z^\prime_v$ whose slope makes the smallest clockwise angle with $\calR$, and the zigzag path $z^{\prime\prime}_v$ whose slope makes the smallest anti-clockwise angle with~$\calR$.
Since $\Gamma$ is properly ordered, these zigzag paths are consecutive around~$v$.
The intersection of~$z^\prime_v$ and $z^{\prime\prime}_v$ is the edge~$e(v)$ which has~$v$ as an endpoint.

\begin{center}
{\scalebox{0.7}{
\begin{tikzpicture}
\newcommand{\edgewidth}{0.05cm} %line_width_of _edges
\newcommand{\nodewidth}{0.05cm} %line_width_around_nodes
\newcommand{\noderad}{0.16} %radius_of_node

%coordinate setting
\coordinate (W1) at (0,0); \coordinate (B1) at (3,0);
\path (W1) ++(135:1.5cm) coordinate (W1a); \path (W1) ++(225:1.5cm) coordinate (W1b);
\path (B1) ++(45:1.5cm) coordinate (B1a); \path (B1) ++(315:1.5cm) coordinate (B1b);

%edge
\draw [line width=\edgewidth] (W1)--(B1); \draw [line width=\edgewidth] (W1)--(W1a); \draw [line width=\edgewidth] (W1)--(W1b);
\draw [line width=\edgewidth] (B1)--(B1a); \draw [line width=\edgewidth] (B1)--(B1b);
\draw [line width=\edgewidth,line cap=round, dash pattern=on 0pt off 2.5\pgflinewidth] (-1,0.35)--(-1,-0.35);
\draw [line width=\edgewidth,line cap=round, dash pattern=on 0pt off 2.5\pgflinewidth] (4,0.35)--(4,-0.35);
\draw [line width=\nodewidth, fill=white] (W1) circle [radius=\noderad] ;
\draw [line width=\nodewidth, fill=black] (B1) circle [radius=\noderad] ;

%zigzag
\path (W1) ++(270:0.3cm) coordinate (W1-); \path (W1) ++(90:0.3cm) coordinate (W1+);
\path (B1) ++(270:0.3cm) coordinate (B1-); \path (B1) ++(90:0.3cm) coordinate (B1+);
\path (W1-) ++(225:1.7cm) coordinate (W1--); \path (W1+) ++(135:1.7cm) coordinate (W1++);
\path (B1-) ++(315:1.7cm) coordinate (B1--); \path (B1+) ++(45:1.7cm) coordinate (B1++);
\draw [->, rounded corners, line width=0.08cm, red] (W1--)--(W1-)--(B1+)--(B1++) ;
\draw [->, rounded corners, line width=0.08cm, blue] (B1--)--(B1-)--(W1+)--(W1++) ;

%node
\node at (-0.5,0) {\large$v$} ; \node at (3.5,0.1) {\large$v^\prime$} ;
\node at (1.5,-0.7) {\large$e(v)=e(v^\prime)$} ;
\node[red] at (4.5,1.7) {\large$z_v^\prime$}; \node[blue] at (-1.5,1.7) {\large$z_v^{\prime\prime}$};

\end{tikzpicture}
} }
\end{center}

We then apply the same argument to the node $v^\prime$ which is the other endpoint of $e(v)$.
The proper ordering on $\Gamma$ leads to the conclusion that $e(v)=e(v^\prime)$ (see the above figure).
We repeat these arguments for any node, but it is enough to consider $e(v)$'s for any $v\in\Gamma^+_0$ (or $v\in\Gamma^-_0$).
By \cite[Section~3.2]{Gul} or \cite[Lemma~4.19]{Bro}, we see that the subset of edges $e(v)$ for all $v\in\Gamma^+_0$ forms a perfect matching.
Furthermore, since $z=z_k$ is type I, there exists a zigzag path whose slope is located at an angle less than $\pi$ in an anti-clockwise (resp.\ clockwise) direction from~$[z]$ in~$\xi(v)$ by \cite[Lemma~4.11 and its proof]{Bro}, in which case such a slope is located between $[z]$ and $[z^\prime_v]$ (resp.~$[z^{\prime\prime}_v]$) or coincides with $[z^\prime_v]$ (resp.~$[z^{\prime\prime}_v]$). Thus, we have $z\neq z^\prime_v$ and $z\neq z^{\prime\prime}_v$ by the definition of the ray $\calR$.
Therefore, in this case $e(v)$ is not contained in~$z$ by the above construction.
Also, we see that if a node~$v$ does not lie on~$z$, $e(v)$ is not contained in~$z$.
Thus, the perfect matching constructed in the above fashion satisfies the desired condition.
\end{proof}

\begin{Lemma}\label{lem_same_length}
Let $z$ be a type I zigzag path on a consistent dimer model.
Then, we have $\omega_{\min}(z)=0$, and hence $\ell(z)$ is the same for all type I zigzag paths having the same slope.
\end{Lemma}

\begin{proof}
This follows from Lemmas~\ref{zigzag_lem2} and \ref{lem_existence_pm}.
\end{proof}

For zigzag paths $z$, $w$ on a dimer model $\Gamma$, we denote by $z\cap w$ the subset of edges that are intersections of $z$ and $w$ on $\Gamma$.
We remark that if $z$ is type I then the number of intersections of~$\widetilde{z}$ and~$\widetilde{w}$ on $\widetilde{\Gamma}$ is less than or equal to one, but there are possibly more intersections of~$z$ and~$w$ if we consider them on~$\Gamma$.

\begin{Lemma}\label{intersect_zigorzag}
Let $z$ be a type I zigzag path on a reduced consistent dimer model $\Gamma$.
We suppose that a zigzag path $w$ has intersections with~$z$ on~$\Gamma$.
Then, we see that $z\cap w\subset \Zig(z)$ or $z\cap w\subset \Zag(z)$.
\end{Lemma}

\begin{proof}
Let $e_1$, $e_2$ be edges of $\Gamma$, and we assume that $w$ intersects with $z$ at $e_1$ and~$e_2$.
If~$e_i$ is a~zig of~$z$, then it is a~zag of~$w$, and vice versa. We assume that $e_1$ is a zig of~$z$ and~$e_2$ is a~zag of~$z$.

We then lift these edges on the universal cover $\widetilde{\Gamma}$. Let $\widetilde{e_1}$, $\widetilde{e_2}$ be edges of $\widetilde{\Gamma}$ whose restrictions on $\Gamma$ are $e_1$, $e_2$, respectively. In particular, $\widetilde{e_1}$ is a~zig of~$\widetilde{z}$ and~$\widetilde{e_2}$ is a~zag of~$\widetilde{z}$.
Then, there exist zigzag paths $\widetilde{w}_1$, $\widetilde{w}_2$ on~$\widetilde{\Gamma}$ whose restriction on~$\Gamma$ is just~$w$,
and $\widetilde{w}_1$ (resp.~$\widetilde{w}_2$) intersects with~$\widetilde{z}$ at~$\widetilde{e_1}$ (resp.~$\widetilde{e_2}$).
The zigzag path $\widetilde{z}$ splits $\RR^2$ into two pieces, and~$\widetilde{w}_1$ (resp.~$\widetilde{w}_2$) intersects with~$\widetilde{z}$ from left to right (resp.\ right to left) by the definition of zigzag paths (see the figure below).

\begin{center}
\newcommand{\edgewidth}{0.05cm} %line_width_of _edges
\newcommand{\nodewidth}{0.05cm} %line_width_around_nodes
\newcommand{\noderad}{0.16} %radius_of_node
\scalebox{0.6}{
\begin{tikzpicture}
%coordinate setting
\coordinate (B1) at (1,2); \coordinate (B2) at (3,2); \coordinate (B3) at (3,4);
\coordinate (B4) at (6,2); \coordinate (B5) at (6,4); \coordinate (B6) at (8,2);

\coordinate (W1) at (2,1); \coordinate (W2) at (2,3); \coordinate (W3) at (4,3);
\coordinate (W4) at (7,1); \coordinate (W5) at (7,3);
%edge
\draw [line width=\edgewidth] (2.5,0.5)--(W1)--(B2)--(W2)--(B3)--(2.5,4.5) ;
\draw [line width=\edgewidth,line cap=round, dash pattern=on 0pt off 2.5\pgflinewidth] (2.6,0.4)--(3,0);
\draw [line width=\edgewidth,line cap=round, dash pattern=on 0pt off 2.5\pgflinewidth] (2.4,4.6)--(2,5);
\draw [line width=\edgewidth] (0.5,2.5)--(B1)--(W2)--(B2)--(W3)--(4.5,2.75) ;
\draw [line width=\edgewidth,line cap=round, dash pattern=on 0pt off 2.5\pgflinewidth] (0.4,2.6)--(0,3);
\draw [line width=\edgewidth,line cap=round, dash pattern=on 0pt off 2.5\pgflinewidth] (4.7,2.65)--(5.4,2.3);
\draw [line width=\edgewidth] (5.5,2.25)--(B4)--(W5)--(B6)--(8.5,2.5) ;
\draw [line width=\edgewidth,line cap=round, dash pattern=on 0pt off 2.5\pgflinewidth] (8.6,2.6)--(9,3);
\draw [line width=\edgewidth] (6.5,0.5)--(W4)--(B4)--(W5)--(B5)--(6.5,4.5) ;
\draw [line width=\edgewidth,line cap=round, dash pattern=on 0pt off 2.5\pgflinewidth] (6.4,0.4)--(6,0);
\draw [line width=\edgewidth,line cap=round, dash pattern=on 0pt off 2.5\pgflinewidth] (6.6,4.6)--(7,5);

%blacknode
\draw [line width=\nodewidth, fill=black] (B1) circle [radius=\noderad] ; \draw [line width=\nodewidth, fill=black] (B2) circle [radius=\noderad] ;
\draw [line width=\nodewidth, fill=black] (B3) circle [radius=\noderad] ;
\draw [line width=\nodewidth, fill=black] (B4) circle [radius=\noderad] ; \draw [line width=\nodewidth, fill=black] (B5) circle [radius=\noderad] ;
\draw [line width=\nodewidth, fill=black] (B6) circle [radius=\noderad] ;
%whitenode
\draw [line width=\nodewidth, fill=white] (W1) circle [radius=\noderad] ; \draw [line width=\nodewidth, fill=white] (W2) circle [radius=\noderad] ;
\draw [line width=\nodewidth, fill=white] (W3) circle [radius=\noderad] ;
\draw [line width=\nodewidth, fill=white] (W4) circle [radius=\noderad] ; \draw [line width=\nodewidth, fill=white] (W5) circle [radius=\noderad] ;

%zigzag
\draw [->, rounded corners, line width=0.08cm, blue] (0,3.25)--(1,2.25)--(2,3.25)--(3,2.25)--(4,3.25)--(6,2.25)--(7,3.25)--(8,2.25)--(9,3.25);
\node[blue] at (9.5,3.25) {\Large$\widetilde{z}$} ;
\draw [->, rounded corners, line width=0.08cm, red] (2.75,0)--(1.75,1)--(2.75,2)--(1.75,3)--(2.75,4)--(1.75,5);
\node[red] at (1.75,5.5) {\Large$\widetilde{w}_2$} ;
\draw [<-, rounded corners, line width=0.08cm, red] (6.25,0)--(7.25,1)--(6.25,2)--(7.25,3)--(6.25,4)--(7.25,5);
\node[red] at (6.25,-0.5) {\Large$\widetilde{w}_1$} ;

\node at (4.5,1) {\Large right of $\widetilde{z}$} ; \node at (4.5,4.5) {\Large left of $\widetilde{z}$} ;
\end{tikzpicture}
}
\end{center}

Since $z$ is type I, there is no intersections of $\widetilde{z}$ and $\widetilde{w}_i$ except $\widetilde{e_i}$ where $i=1,2$.
Thus, we can not superimpose $\widetilde{w}_1$ and $\widetilde{w}_2$ using translations.
This contradicts a choice of $\widetilde{w}_1$, $\widetilde{w}_2$.
\end{proof}

The following lemma follows from the argument in the proof of \cite[Proposition~3.12]{Bro}.

\begin{Lemma}\label{slope_linearly_independent}
Let $z$, $w$ be zigzag paths on a consistent dimer model~$\Gamma$.
We assume that $\widetilde{z}$ intersects with $\widetilde{w}$ on the universal cover $\widetilde{\Gamma}$ at most once.
Then, the slopes $[z]$, $[w]$ are linearly independent if and only if~$\widetilde{z}$ and~$\widetilde{w}$ intersect on precisely one edge.
\end{Lemma}

\begin{Lemma}\label{intersect_zigorzag2}
Let $z_1,\dots,z_r$ be type~I zigzag paths on a reduced consistent dimer model $\Gamma$ having the same slope.
We suppose that a zigzag path~$w$ has intersections with $z_j$ for some $j$.
Then, $w$~intersects with every $z_i$ $(i=1,\dots,r)$, and the intersections $w\cap z_1,\dots, w\cap z_r$ are all in $\Zig(w)$ or all in $\Zag(w)$.
Moreover, we have $|w\cap z_1|=\cdots=|w\cap z_r|$.
\end{Lemma}

\begin{proof}
Since $z_1,\dots,z_r$ are type I, $\widetilde{w}$ intersects with each $\widetilde{z_i}$ at most once.
Thus, each pair of zigzag paths $(\widetilde{w},\widetilde{z_i})$ for $i=1,\dots,r$ satisfies the assumption in Lemma~\ref{slope_linearly_independent}.
Since~$w$ has intersections with~$z_j$, $\widetilde{w}$ intersects with $\widetilde{z_j}$ precisely once on the universal cover.
By Lemma~\ref{slope_linearly_independent}, we see that $[w]$ and $[z_j]$ are linearly independent, and hence $[w]$ and $[z_i]$ are linearly independent for any~$i$.
Thus, $\widetilde{w}$ intersects with $\widetilde{z_i}$ precisely once for any $i=1,\dots,r$.

The latter assertion follows from a similar argument as in the proof of Lemma~\ref{intersect_zigorzag}.
More precisely, since $\widetilde{w}$ intersects with $\widetilde{z_i}$ precisely once for any $i=1,\dots,r$,
if $\widetilde{w}$ intersects with $\widetilde{z_i}$ from right to left (resp.\ left to right), then so do zigzag paths having the same slope.
This means that intersections are all in $\Zig(w)$ (resp.\ all in $\Zag(w))$.

Also, let $\widetilde{z_1}$ and $\widetilde{z_1}^\prime$ be zigzag paths on $\widetilde{\Gamma}$ that are projected onto the zigzag path $z_1$ on $\Gamma$.
We assume that there is no zigzag path projected onto $z_1$ between $\widetilde{z_1}$ and $\widetilde{z_1}^\prime$.
Also, we assume that~$\widetilde{w}$ first intersects with~$\widetilde{z_1}$, then intersects with $\widetilde{z_1}^\prime$.
Then, for all $i=2,\dots,r$ we can find a~unique zigzag path~$\widetilde{z_i}$ on $\widetilde{\Gamma}$ such that it is projected onto $z_i$ and is located between~$\widetilde{z_1}$ and~$\widetilde{z_1}^\prime$.
Thus, after $\widetilde{w}$ intersects with $\widetilde{z_1}$, it intersects with $\widetilde{z_2}, \dots, \widetilde{z_r}$ precisely once and then arrives at~$\widetilde{z_1}^\prime$.
We can apply the same arguments for any pair $(\widetilde{z_1},\widetilde{z_1}^\prime)$ of zigzag paths on~$\widetilde{\Gamma}$ satisfying the above properties. Thus, projecting onto~$\Gamma$ we have $|w\cap z_1|=\cdots=|w\cap z_r|$.
\end{proof}
